% Section 9 -- Conclusions. PR: longer than the abstract; strengths, weaknesses, benefit to others, future work.
\section{Conclusions}\label{sec:conclusions}

When each site contributes a single diagnosis, subject-level cross-validation rewards models that score by site, and without an assumption that the data cannot test, diagnosis and site cannot be separated at all. Holding out whole sites removed an optimism of $0.37$ to $0.46$ in AUC; with bridging subjects, what mattered was that they crossed the class boundary, and $L$ ranked designs before training; training only on within-site comparisons kept the reverse-set AUC above chance, with test AUC within $0.01$ of the pooled model when every site was bridged. When every site held both diagnoses in unequal proportions, the shortcut and the optimism of subject-level cross-validation grew with $L$, and within-site training again removed the shortcut.

The main strength of the framework is that it acts on the study design and the training objective, so it can be planned before data collection and combined with standard classifiers, and each rule was tested with pre-specified criteria on three cohorts covering two disorders that share no participants. Its main weaknesses are that class purity was constructed rather than observed and that the theory covers additive site effects and linear scores (Section~\ref{sec:limitations}).

For researchers who pool clinical data, Table~\ref{tab:recommend} lists what to do in each situation; the two simplest steps, tabulating subjects by site and diagnosis and reporting site-level performance, apply to every study, and a few bridging participants per site are worth more than as many additions of each site's own diagnosis; with within-site training two per site were enough, whereas pooled training needed bridges to make up a third of each site. For spurious correlations without bias-conflicting samples in pattern recognition, our results point to a remedy at the level of the data: collecting a few such samples by design. Future work should recruit bridging subjects prospectively at class-pure clinical sites, extend the theory to site-by-diagnosis interactions and non-linear scores, and test within-site objectives for EEG, structural MRI and end-to-end deep models.
